\chapter[Présence, Paie et Analytics]{Gestion de la Présence, Paie et Analyse}
\label{chap:presence_salary_analytics}
\chapterplan{%
\chapterplanitem{}{Introduction}{sec:chap6_introduction}
\chapterplanitem{1}{Sprint 5 : Gestion de présence et paie}{sec:sprint5}
\chapterplanitem{2}{Sprint 6 : Analytics, CI/CD et déploiement}{sec:sprint6}
\chapterplanitem{}{Conclusion}{sec:chap6_conclusion}
}

\section*{Introduction}
\phantomsection
\label{sec:chap6_introduction}

Ce chapitre présente la troisième et dernière release fonctionnelle du projet. Après le socle
d'authentification et de gestion des leads (Release 1) et le moteur d'IA avec la gestion de la
qualité (Release 2), cette release finalise la plateforme CRM en couvrant les aspects
opérationnels et analytiques.

Le Sprint 5 traite de la gestion de présence (pointage, pauses) et du calcul automatisé des
salaires. Le Sprint 6 complète le projet avec les tableaux de bord analytics, l'export de
rapports, l'intégration continue et le déploiement Docker.

%% ══════════════════════════════════════════════════════════════════════════════
\section[Sprint 5 : Présence et Paie]{Sprint 5 : Gestion de présence et paie}
\label{sec:sprint5}

%% ────────────────────────────────────────────────────────────────────────────
\subsection{Objectifs et Backlog du Sprint 5}
\label{subsec:sprint5_backlog}

Le Sprint 5 a pour objectif de livrer les fonctionnalités de gestion des ressources humaines
de la plateforme~: le système de pointage (clock-in/clock-out), la gestion des pauses, le
suivi de présence en temps réel et le calcul automatisé des salaires.

\begingroup
\small
\setlength{\tabcolsep}{5pt}
\renewcommand{\arraystretch}{1.35}
\begin{longtable}{|>{\centering\arraybackslash}p{1.1cm}|>{\raggedright\arraybackslash}p{9.0cm}|>{\centering\arraybackslash}p{1.6cm}|>{\centering\arraybackslash}p{1.6cm}|}
\caption{Backlog du Sprint 5 -- Gestion de présence et paie}\label{tab:sprint5_backlog}\\
\hline
\textbf{ID} & \textbf{User Story} & \textbf{Priorité} & \textbf{Statut} \\
\hline
\endfirsthead
\hline
\textbf{ID} & \textbf{User Story} & \textbf{Priorité} & \textbf{Statut} \\
\hline
\endhead
\hline
\endfoot
\endlastfoot
6.1 & En tant qu'\textbf{Agent}, je veux pointer mon arrivée (clock-in) afin de déclarer ma présence. & Élevée & Terminé \\
\hline
6.2 & En tant qu'\textbf{Agent}, je veux pointer mon départ (clock-out) afin de clôturer ma journée. & Élevée & Terminé \\
\hline
6.3 & En tant qu'\textbf{Agent}, je veux déclarer une pause avec son type afin de tracer mes absences. & Moyenne & Terminé \\
\hline
6.4 & En tant qu'\textbf{Agent}, je veux terminer ma pause afin de reprendre mon activité. & Moyenne & Terminé \\
\hline
6.5 & En tant qu'\textbf{Admin}, je veux consulter l'état de présence en temps réel de toute l'équipe. & Élevée & Terminé \\
\hline
6.6 & En tant qu'\textbf{Admin}, je veux consulter le rapport de présence par agent et par période. & Élevée & Terminé \\
\hline
6.7 & En tant qu'\textbf{Admin}, je veux configurer les règles de calcul des salaires (primes, pénalités). & Élevée & Terminé \\
\hline
6.8 & En tant que \textbf{système}, je veux calculer automatiquement les salaires selon les règles configurées. & Élevée & Terminé \\
\hline
6.9 & En tant qu'\textbf{Admin}, je veux consulter le récapitulatif des salaires par mois. & Élevée & Terminé \\
\hline
6.10 & En tant qu'\textbf{Agent}, je veux consulter mon propre salaire et son détail. & Moyenne & Terminé \\
\hline
\end{longtable}
\endgroup

%% ────────────────────────────────────────────────────────────────────────────
\subsection{Spécification et Conception du Sprint 5}
\label{subsec:sprint5_spec}

\subsubsection{Système de pointage}

Le système de pointage permet à chaque agent de gérer sa présence via des actions
élémentaires~:

\begin{itemize}
    \item \textbf{Clock-in}~: enregistre l'heure d'arrivée de l'agent avec le statut
    \texttt{active}. Une seule session active par agent à la fois.
    \item \textbf{Clock-out}~: enregistre l'heure de départ et calcule la durée totale
    de présence, pauses déduites.
    \item \textbf{Pause}~: permet à l'agent de déclarer une pause (déjeuner, café,
    technique) avec des types prédéfinis. Chaque pause est tracée dans la table
    \texttt{AttendanceBreak}.
    \item \textbf{Statut en temps réel}~: expose l'état actuel de chaque agent
    (présent, en pause, absent) pour le suivi par l'Admin.
\end{itemize}

\subsubsection{Modèle de données de présence}

Le diagramme ci-dessous illustre les entités de présence~:

\begin{figure}[H]
    \centering
    \fbox{\parbox{0.85\linewidth}{
    \centering\vspace{0.4cm}
    \textit{[Diagramme de classes : Gestion de présence]}\\[0.3cm]
    Entités impliquées~: User (agent) 1---* Attendance (clockIn, clockOut, status, date),
    Attendance 1---* AttendanceBreak (breakStart, breakEnd, type). Un agent peut avoir
    plusieurs pointages sur différentes journées, chaque pointage pouvant contenir
    plusieurs pauses.
    \vspace{0.4cm}
    }}
    \caption{Modèle de données de la gestion de présence}
    \label{fig:sprint5_attendance_model}
\end{figure}

\subsubsection{Calcul automatisé des salaires}

Le calcul des salaires repose sur un système de règles configurables (\texttt{SalaryRule})
permettant de définir~:

\begin{itemize}
    \item \textbf{Salaire de base}~: montant fixe par agent.
    \item \textbf{Primes sur rendez-vous}~: bonus par RDV réalisé (\texttt{rdv\_bonus}).
    \item \textbf{Primes sur pose}~: bonus par installation (\texttt{pose\_bonus}).
    \item \textbf{Prime qualité}~: bonus basé sur le taux de qualité (\texttt{quality\_rate}).
    \item \textbf{Pénalités}~: déductions pour absences ou non-respect des objectifs.
\end{itemize}

Le calcul est déclenché manuellement par l'Admin via l'API \texttt{GET /api/salaries/calculate}
et produit un résultat détaillé pour chaque agent, incluant le décompte des RDV, des poses,
des refus, le taux de qualité et le salaire total.

%% ────────────────────────────────────────────────────────────────────────────
\subsection{Présentation des Technologies Utilisées}
\label{subsec:sprint5_tech}

\subsubsection{Service de présence}

Le service \texttt{AttendanceService} gère l'ensemble des opérations de pointage~:

\begin{itemize}
    \item \textbf{ClockInAsync}~: crée une nouvelle session de présence pour l'agent.
    \item \textbf{ClockOutAsync}~: ferme la session active et calcule la durée.
    \item \textbf{StartBreakAsync / EndBreakAsync}~: gère le cycle des pauses.
    \item \textbf{GetTeamStatusAsync}~: retourne l'état en temps réel de tous les agents.
    \item \textbf{GetTeamReportAsync}~: génère un rapport de présence pour la période.
\end{itemize}

\subsubsection{Service de paie}

Le service \texttt{SalaryService} implémente les algorithmes de calcul~:

\begin{itemize}
    \item \textbf{CalculateAllSalariesAsync}~: parcourt tous les agents, récupère les
    statistiques du mois (RDV, poses, refus, qualité) et applique les règles.
    \item \textbf{GetSalariesAsync}~: retourne la liste des salaires calculés pour un mois.
    \item \textbf{GetSalaryRulesAsync}~: retourne les règles de calcul configurables.
    \item \textbf{CreateSalaryRuleAsync / UpdateSalaryRuleAsync}~: gestion des règles.
\end{itemize}

%% ────────────────────────────────────────────────────────────────────────────
\subsection{Réalisation}
\label{subsec:sprint5_realisation}

\subsubsection{Interface de pointage pour l'agent}

L'agent dispose d'une interface simple pour gérer sa présence quotidienne.

\begin{figure}[H]
    \centering
    \fbox{\parbox{0.9\linewidth}{
    \centering\vspace{0.4cm}
    \textit{[Capture d'écran : Interface de pointage]}\\[0.3cm]
    Interface de pointage pour l'agent avec boutons Clock-in / Clock-out, statut actuel
    (présent, en pause, absent), chronomètre de la journée en cours et historique des
    pointages récents.
    \vspace{0.4cm}
    }}
    \caption{Interface de pointage agent}
    \label{fig:sprint5_clock_interface}
\end{figure}

\subsubsection{Tableau de bord de présence pour l'Admin}

L'Admin visualise l'état de présence de toute l'équipe en temps réel.

\begin{figure}[H]
    \centering
    \fbox{\parbox{0.9\linewidth}{
    \centering\vspace{0.4cm}
    \textit{[Capture d'écran : Dashboard présence équipe]}\\[0.3cm]
    Tableau de bord présentant la liste des agents avec leur statut actuel
    (présent/absent/en pause), heures d'arrivée, durée de présence, pauses en cours,
    et un récapitulatif hebdomadaire.
    \vspace{0.4cm}
    }}
    \caption{Dashboard de présence équipe}
    \label{fig:sprint5_team_dashboard}
\end{figure}

\subsubsection{Configuration des règles de salaire}

L'Admin peut configurer les règles de calcul des salaires via une interface dédiée.

\begin{figure}[H]
    \centering
    \fbox{\parbox{0.9\linewidth}{
    \centering\vspace{0.4cm}
    \textit{[Capture d'écran : Configuration des règles salariales]}\\[0.3cm]
    Interface listant les règles de salaire avec leur type (rdv\_bonus, pose\_bonus,
    quality\_bonus, penalty), leur montant, le rôle ciblé et leur état actif/inactif.
    Possibilité d'ajouter, modifier ou désactiver une règle.
    \vspace{0.4cm}
    }}
    \caption{Configuration des règles salariales}
    \label{fig:sprint5_salary_rules}
\end{figure}

\subsubsection{Calcul et consultation des salaires}

\begin{figure}[H]
    \centering
    \fbox{\parbox{0.9\linewidth}{
    \centering\vspace{0.4cm}
    \textit{[Capture d'écran : Récapitulatif des salaires]}\\[0.3cm]
    Tableau récapitulatif des salaires par agent pour un mois donné~: salaire de base,
    nombre de RDV, bonus RDV, nombre de poses, bonus pose, taux qualité, prime qualité,
    pénalités et salaire total. Possibilité de filtrer par mois et d'exporter.
    \vspace{0.4cm}
    }}
    \caption{Récapitulatif des salaires}
    \label{fig:sprint5_salary_summary}
\end{figure}

\subsubsection{Tests unitaires et validation}

Les tests du Sprint 5 couvrent~: le service de pointage (clock-in/out, chevauchement,
pauses), le calcul des salaires (règles, bonus, pénalités), les rapports de présence,
les scénarios d'erreur (clock-out sans clock-in, pause sans session active).

\subsubsection{Résultats obtenus}

Le Sprint 5 livre un système complet de gestion des ressources humaines. Les agents
peuvent gérer leur présence de manière autonome. L'Admin dispose d'outils de suivi
en temps réel et de rapports de présence. Le calcul automatisé des salaires, basé
sur des règles configurables, réduit considérablement le temps de traitement de la
paie et élimine les erreurs manuelles.

%% ══════════════════════════════════════════════════════════════════════════════
\section[Sprint 6 : Analytics et CI/CD]{Sprint 6 : Analytics, CI/CD et déploiement}
\label{sec:sprint6}

%% ────────────────────────────────────────────────────────────────────────────
\subsection{Objectifs et Backlog du Sprint 6}
\label{subsec:sprint6_backlog}

Le Sprint 6 finalise la plateforme avec les tableaux de bord analytics, les exports
de rapports, le pipeline d'intégration continue et le déploiement conteneurisé.

\begingroup
\small
\setlength{\tabcolsep}{5pt}
\renewcommand{\arraystretch}{1.35}
\begin{longtable}{|>{\centering\arraybackslash}p{1.1cm}|>{\raggedright\arraybackslash}p{9.0cm}|>{\centering\arraybackslash}p{1.6cm}|>{\centering\arraybackslash}p{1.6cm}|}
\caption{Backlog du Sprint 6 -- Analytics et CI/CD}\label{tab:sprint6_backlog}\\
\hline
\textbf{ID} & \textbf{User Story} & \textbf{Priorité} & \textbf{Statut} \\
\hline
\endfirsthead
\hline
\textbf{ID} & \textbf{User Story} & \textbf{Priorité} & \textbf{Statut} \\
\hline
\endhead
\hline
\endfoot
\endlastfoot
7.1 & En tant qu'\textbf{Admin}, je veux consulter un dashboard analytics global avec les KPIs du centre d'appels. & Élevée & Terminé \\
\hline
7.2 & En tant qu'\textbf{Admin}, je veux visualiser la performance individuelle de chaque agent. & Élevée & Terminé \\
\hline
7.3 & En tant qu'\textbf{Admin}, je veux consulter les statistiques géographiques des leads. & Moyenne & Terminé \\
\hline
7.4 & En tant qu'\textbf{Admin}, je veux suivre les suivis (follow-ups) par agent. & Moyenne & Terminé \\
\hline
7.5 & En tant qu'\textbf{Admin}, je veux surveiller l'activité en temps réel du centre d'appels. & Élevée & Terminé \\
\hline
7.6 & En tant qu'\textbf{Admin}, je veux exporter les rapports en PDF/Excel afin de les diffuser. & Moyenne & Terminé \\
\hline
7.7 & En tant que \textbf{système}, je veux exécuter les tests automatiquement à chaque push via GitHub Actions. & Élevée & Terminé \\
\hline
7.8 & En tant que \textbf{système}, je veux être déployé via Docker Compose avec trois services (API, Frontend, PostgreSQL). & Élevée & Terminé \\
\hline
\end{longtable}
\endgroup

%% ────────────────────────────────────────────────────────────────────────────
\subsection{Spécification et Conception du Sprint 6}
\label{subsec:sprint6_spec}

\subsubsection{Tableaux de bord analytics}

La plateforme propose plusieurs vues analytiques~:

\begin{itemize}
    \item \textbf{Dashboard global}~: indicateurs clés (nombre d'appels, score moyen,
    taux de conversion, agents actifs, leads par statut).
    \item \textbf{Performance agents}~: classement des agents par score, taux de
    conversion, nombre de RDV, durée d'appel moyenne et taux de présence.
    \item \textbf{Supervision}~: vue d'ensemble des métriques opérationnelles avec
    graphiques d'évolution temporelle.
    \item \textbf{Analyse géographique}~: répartition des leads par code postal et
    par région.
    \item \textbf{Temps réel}~: activité en direct du centre d'appels (appels en
    cours, agents en ligne, files d'attente).
\end{itemize}

\subsubsection{Pipeline CI/CD}

Le pipeline d'intégration et de déploiement continu est orchestré via GitHub Actions~:

\begin{itemize}
    \item \textbf{Déclencheur}~: push sur les branches \texttt{main} et \texttt{develop}.
    \item \textbf{Backend}~: restauration des packages NuGet, build, exécution des tests
    unitaires xUnit, publication du build.
    \item \textbf{Frontend}~: installation des dépendances, build TypeScript, exécution
    des tests Vitest, linting.
    \item \textbf{Déploiement}~: construction des images Docker et push vers le registry.
\end{itemize}

\subsubsection{Infrastructure Docker}

L'infrastructure de déploiement est définie via \texttt{docker-compose.yml} avec trois
services~:

\begin{itemize}
    \item \textbf{crm-api}~: application .NET 8 avec ASP.NET Core, expose le port 5000.
    \item \textbf{crm-frontend}~: application React 18 servie via Nginx, expose le port 80.
    \item \textbf{crm-db}~: base de données PostgreSQL 16, volume persistant, expose
    le port 5432.
\end{itemize}

Le déploiement inclut des healthchecks pour chaque service, une configuration de
redémarrage automatique et un réseau interne pour la communication inter-services.

%% ────────────────────────────────────────────────────────────────────────────
\subsection{Présentation des Technologies Utilisées}
\label{subsec:sprint6_tech}

\subsubsection{AnalyticsService}

Le service \texttt{AnalyticsService} centralise les requêtes analytiques~:

\begin{itemize}
    \item \textbf{GetOverviewAsync}~: KPIs globaux (appels, scores, agents, leads, RDV).
    \item \textbf{GetAgentsPerformanceAsync}~: performance détaillée par agent.
    \item \textbf{GetSupervisionAsync}~: métriques de supervision et tendances.
    \item \textbf{GetGeoAsync}~: données géographiques des leads.
    \item \textbf{GetFollowupsAsync}~: statistiques de suivi par agent.
    \item \textbf{GetLiveAgentsAsync}~: état en temps réel des agents.
\end{itemize}

\subsubsection{GitHub Actions}

Le pipeline CI/CD utilise GitHub Actions avec les étapes suivantes~:

\begin{itemize}
    \item \textbf{Checkout}~: récupération du code source.
    \item \textbf{Setup .NET 8}~: installation du SDK .NET 8.
    \item \textbf{Restore \& Build}~: restauration des dépendances et compilation.
    \item \textbf{Test}~: exécution des tests xUnit avec couverture de code.
    \item \textbf{Setup Node.js}~: installation de Node.js pour le frontend.
    \item \textbf{Frontend Build \& Test}~: build et tests Vitest du frontend.
    \item \textbf{Docker Build}~: construction et push des images Docker.
\end{itemize}

\subsubsection{Docker et Nginx}

Le frontend React est servi via Nginx en mode production. Une configuration
\texttt{nginx.conf} gère le routage SPA, la compression gzip et les proxys
vers l'API backend pour les requêtes \texttt{/api/*}.

%% ────────────────────────────────────────────────────────────────────────────
\subsection{Réalisation}
\label{subsec:sprint6_realisation}

\subsubsection{Dashboard analytics global}

\begin{figure}[H]
    \centering
    \fbox{\parbox{0.9\linewidth}{
    \centering\vspace{0.4cm}
    \textit{[Capture d'écran : Dashboard analytics global]}\\[0.3cm]
    Dashboard présentant les KPIs principaux~: nombre total d'appels, score moyen,
    taux de conversion, agents actifs, leads par statut (graphique circulaire),
    évolution des appels sur la période (graphique en aires) et classement des agents.
    \vspace{0.4cm}
    }}
    \caption{Dashboard analytics global}
    \label{fig:sprint6_analytics_dashboard}
\end{figure}

\subsubsection{Performance des agents}

\begin{figure}[H]
    \centering
    \fbox{\parbox{0.9\linewidth}{
    \centering\vspace{0.4cm}
    \textit{[Capture d'écran : Performance des agents]}\\[0.3cm]
    Tableau de bord de performance individuelle~: classement des agents par score
    moyen, taux de conversion, durée moyenne d'appel, ratio de parole agent/client,
    nombre de RDV obtenus et tendance d'évolution.
    \vspace{0.4cm}
    }}
    \caption{Performance des agents}
    \label{fig:sprint6_agent_performance}
\end{figure}

\subsubsection{Activité en temps réel}

\begin{figure}[H]
    \centering
    \fbox{\parbox{0.9\linewidth}{
    \centering\vspace{0.4cm}
    \textit{[Capture d'écran : Vue temps réel]}\\[0.3cm]
    Interface de monitoring en temps réel affichant les appels en cours, les agents
    actifs, le temps d'attente moyen, les alertes récentes et un flux d'activité
    mis à jour automatiquement via WebSocket.
    \vspace{0.4cm}
    }}
    \caption{Vue temps réel du centre d'appels}
    \label{fig:sprint6_realtime}
\end{figure}

\subsubsection{Exports et rapports}

La plateforme permet l'export des données analytiques au format PDF et CSV via le
service \texttt{ExportController}, avec la bibliothèque Recharts pour les graphiques
dans les rapports PDF.

\subsubsection{Architecture de déploiement}

\begin{figure}[H]
    \centering
    \fbox{\parbox{0.95\linewidth}{
    \centering\vspace{0.4cm}
    \textit{[Diagramme d'architecture de déploiement]}\\[0.3cm]
    Architecture Docker Compose avec trois conteneurs~: Nginx (port 80) sert le
    frontend React 18 et proxy les requêtes /api/* vers l'API .NET 8 (port 5000),
    qui communique avec PostgreSQL 16 (port 5432). Les healthchecks assurent la
    disponibilité de chaque service. Le pipeline GitHub Actions construit et
    déploie automatiquement sur push vers main.
    \vspace{0.4cm}
    }}
    \caption{Architecture de déploiement Docker}
    \label{fig:sprint6_deployment_architecture}
\end{figure}

\subsubsection{Tests et validation}

Les tests du Sprint 6 couvrent~: le service analytics (KPIs, performances, calculs
statistiques), les exports, les scénarios de déploiement Docker, la configuration
Nginx et la validation du pipeline CI/CD.

\subsubsection{Résultats obtenus}

Le Sprint 6 finalise la plateforme CRM intelligente. Les tableaux de bord analytics
offrent une vision complète de l'activité du centre d'appels, permettant aux managers
de piloter la performance en temps réel. Le pipeline CI/CD garantit la qualité du
code à chaque livraison, et l'infrastructure Docker assure un déploiement
reproductible et scalable.

%% ══════════════════════════════════════════════════════════════════════════════
\section*{Conclusion}
\phantomsection
\label{sec:chap6_conclusion}

Ce chapitre a présenté la troisième et dernière release de la plateforme CRM. Le Sprint 5
a livré un système complet de gestion des ressources humaines~: pointage en temps réel,
gestion des pauses et calcul automatisé des salaires basé sur des règles configurables.
Le Sprint 6 a finalisé le projet avec des tableaux de bord analytics riches, l'export
de rapports, un pipeline CI/CD robuste via GitHub Actions et un déploiement
conteneurisé via Docker Compose.

À l'issue de cette release, la plateforme CRM intelligente couvre l'ensemble du cycle
de vie d'un centre d'appels~: de l'acquisition de leads à la rémunération des agents,
en passant par le scoring IA, la gestion de la qualité et l'analyse de performance.
Le tout repose sur une architecture moderne, testée et déployable en environnements
professionnels.
